\documentclass[10pt,a4paper]{amsart}
\usepackage{amsmath,amssymb,amsthm,graphicx}
\usepackage{verbatim}
\newenvironment{reference}{\comment}{\endcomment}
%\newenvironment{reference}{}{}
\newenvironment{slogan}{\comment}{\endcomment}
\newenvironment{history}{\comment}{\endcomment}

% For commutative diagrams we use Xy-pic
\usepackage[all]{xy}

% We use 2cell for 2-commutative diagrams.
\xyoption{2cell}
% Theorem environments.
%
\theoremstyle{plain}
\newtheorem{theorem}[subsection]{Theorem}
\newtheorem{proposition}[subsection]{Proposition}
\newtheorem{lemma}[subsection]{Lemma}

\theoremstyle{definition}
\newtheorem{definition}[subsection]{Definition}
\newtheorem{example}[subsection]{Example}
\newtheorem{exercise}[subsection]{Exercise}
\newtheorem{situation}[subsection]{Situation}

\theoremstyle{remark}
\newtheorem{remark}[subsection]{Remark}
\newtheorem{remarks}[subsection]{Remarks}

\numberwithin{equation}{subsection}

% Macros
%
\def\lim{\mathop{\mathrm{lim}}\nolimits}
\def\colim{\mathop{\mathrm{colim}}\nolimits}
\def\Spec{\mathop{\mathrm{Spec}}}
\def\Hom{\mathop{\mathrm{Hom}}\nolimits}
\def\Ext{\mathop{\mathrm{Ext}}\nolimits}
\def\SheafHom{\mathop{\mathcal{H}\!\mathit{om}}\nolimits}
\def\SheafExt{\mathop{\mathcal{E}\!\mathit{xt}}\nolimits}
\def\Sch{\mathit{Sch}}
\def\Mor{\mathop{\mathrm{Mor}}\nolimits}
\def\Ob{\mathop{\mathrm{Ob}}\nolimits}
\def\Sh{\mathop{\mathit{Sh}}\nolimits}
\def\NL{\mathop{N\!L}\nolimits}
\def\CH{\mathop{\mathrm{CH}}\nolimits}
\def\proetale{{pro\text{-}\acute{e}tale}}
\def\etale{{\acute{e}tale}}
\def\QCoh{\mathit{QCoh}}
\def\Ker{\mathop{\mathrm{Ker}}}
\def\Im{\mathop{\mathrm{Im}}}
\def\Coker{\mathop{\mathrm{Coker}}}
\def\Coim{\mathop{\mathrm{Coim}}}

% Boxtimes
%
\DeclareMathSymbol{\boxtimes}{\mathbin}{AMSa}{"02}

%
% Macros for moduli stacks/spaces
%
\def\QCohstack{\mathcal{QC}\!\mathit{oh}}
\def\Cohstack{\mathcal{C}\!\mathit{oh}}
\def\Spacesstack{\mathcal{S}\!\mathit{paces}}
\def\Quotfunctor{\mathrm{Quot}}
\def\Hilbfunctor{\mathrm{Hilb}}
\def\Curvesstack{\mathcal{C}\!\mathit{urves}}
\def\Polarizedstack{\mathcal{P}\!\mathit{olarized}}
\def\Complexesstack{\mathcal{C}\!\mathit{omplexes}}
% \Pic is the operator that assigns to X its picard group, usage \Pic(X)
% \Picardstack_{X/B} denotes the Picard stack of X over B
% \Picardfunctor_{X/B} denotes the Picard functor of X over B
\def\Pic{\mathop{\mathrm{Pic}}\nolimits}
\def\Picardstack{\mathcal{P}\!\mathit{ic}}
\def\Picardfunctor{\mathrm{Pic}}
\def\Deformationcategory{\mathcal{D}\!\mathit{ef}}
\usepackage{cleveref}
\crefname{originallemma}{original lemma}{original lemmas}
\Crefname{originallemma}{Original lemma}{Original lemmas}
\makeatletter
\newlabel{lemma-base-change-finite-type}{{03XH}{source}{}{sourcepointer-1}{}}
\newlabel{lemma-base-change-finite-type@cref}{{[originallemma][0][]03XH}{[1][0][]source}{}{}{}}
\makeatother
\crefname{originallemma}{original lemma}{original lemmas}
\Crefname{originallemma}{Original lemma}{Original lemmas}
\makeatletter
\newlabel{lemma-composition-finite-type}{{03XG}{source}{}{sourcepointer-2}{}}
\newlabel{lemma-composition-finite-type@cref}{{[originallemma][0][]03XG}{[1][0][]source}{}{}{}}
\makeatother
\crefname{originallemma}{original lemma}{original lemmas}
\Crefname{originallemma}{Original lemma}{Original lemmas}
\makeatletter
\newlabel{spaces-properties-lemma-cover-by-union-affines}{{03FX}{source}{}{sourcepointer-3}{}}
\newlabel{spaces-properties-lemma-cover-by-union-affines@cref}{{[originallemma][0][]03FX}{[1][0][]source}{}{}{}}
\makeatother
\title{Nine equivalent formulations of local finite typeness for morphisms of algebraic spaces}
\author{The Stacks Project Authors}
\date{stacks-project:spaces-morphisms.tex:local-snapshot-20261012; GNU FDL 1.2 or later}
\begin{document}
\maketitle
\setcounter{subsection}{1}
\noindent\textit{Editorial note.} All delivered mathematical text is the author's own, copied from the pinned Stacks Project chapter file spaces-morphisms.tex; no statement or proof step was rewritten, reordered or removed.\par\medskip
\noindent\textit{Editorial note.} The proof states that the implications (2), (3), (4), (5), (6), (7), (9) implies (8) are straightforward and then carries one of them out as the model case; that model case plus the explicit equivalences (1)/(7)/(8) and the final implications cover every arrow of the diagram of conditions.\par\medskip
\noindent\textit{Editorial note.} The three references of the counted spans are recorded as non-load-bearing navigation pointers; the argument that consumes them stays inside the counted proof span.\par\medskip
\noindent\textit{Editorial note.} Preamble: the three native excerpt spans of sources/stacks-spaces-morphisms/preamble.tex (lines 8-13, 14-18 and 154-217). The chapter file itself carries no preamble text.\par\medskip
\noindent\textit{Allowed standard prerequisites (editorial).}\begin{itemize}
\item Covering a scheme by a disjoint union of affines and covering by affine etale neighbourhoods. Properties of Spaces, Lemma spaces-properties-lemma-cover-by-union-affines, used at proof lines 4658-4661. Provides the cover V used to pass from condition (5) to condition (8).
\item Permanence of local finite typeness under composition. Lemma lemma-composition-finite-type, used at proof lines 4663-4665. Makes U -> V locally of finite type in the reduction to (8).
\item Stability of local finite typeness under base change. Lemma lemma-base-change-finite-type, used at proof lines 4674-4676. Gives conditions (3), (4), (5) and (9) from condition (1).
\item The definition of a morphism of finite type at a point and its agreement with the intrinsic condition. used at proof lines 4656-4669 and 4672-4673. Underlies the equivalence of (1), (7) and (8) and the implication to (2).
\end{itemize}\medskip
\ifdefined\hypertarget\hypertarget{sourcepointer-1}{}\fi
\noindent\textit{Source reference.} Original lemma 03XH: focused-seeds/AGP/sources/stacks-snapshot-20261011/tags.txt, source lines 5148--5148. Navigation pointer into the Stacks Project; the standard fact it supplies is applied inside the counted proof span and the argument resting on it stays there.\par\medskip
\ifdefined\hypertarget\hypertarget{sourcepointer-2}{}\fi
\noindent\textit{Source reference.} Original lemma 03XG: focused-seeds/AGP/sources/stacks-snapshot-20261011/tags.txt, source lines 5147--5147. Navigation pointer into the Stacks Project; the standard fact it supplies is applied inside the counted proof span and the argument resting on it stays there.\par\medskip
\ifdefined\hypertarget\hypertarget{sourcepointer-3}{}\fi
\noindent\textit{Source reference.} Original lemma 03FX: focused-seeds/AGP/sources/stacks-snapshot-20261011/tags.txt, source lines 4508--4508. Navigation pointer into the Stacks Project; the standard fact it supplies is applied inside the counted proof span and the argument resting on it stays there.\par\medskip
\begin{lemma}
\label{lemma-finite-type-local}
Let $S$ be a scheme.
Let $f : X \to Y$ be a morphism of algebraic spaces over $S$.
The following are equivalent:
\begin{enumerate}
\item $f$ is locally of finite type,
\item for every $x \in |X|$ the morphism $f$ is of finite type at $x$,
\item for every scheme $Z$ and any morphism $Z \to Y$ the morphism
$Z \times_Y X \to Z$ is locally of finite type,
\item for every affine scheme $Z$ and any morphism
$Z \to Y$ the morphism $Z \times_Y X \to Z$ is locally of finite type,
\item there exists a scheme $V$ and a surjective \'etale morphism
$V \to Y$ such that $V \times_Y X \to V$ is locally of finite type,
\item there exists a scheme $U$ and a surjective \'etale morphism
$\varphi : U \to X$ such that the composition $f \circ \varphi$
is locally of finite type,
\item for every commutative diagram
$$
\xymatrix{
U \ar[d] \ar[r] & V \ar[d] \\
X \ar[r] & Y
}
$$
where $U$, $V$ are schemes and the vertical arrows are \'etale
the top horizontal arrow is locally of finite type,
\item there exists a commutative diagram
$$
\xymatrix{
U \ar[d] \ar[r] & V \ar[d] \\
X \ar[r] & Y
}
$$
where $U$, $V$ are schemes, the vertical arrows are \'etale, $U \to X$
is surjective, and the top horizontal arrow is locally of finite type, and
\item there exist Zariski coverings $Y = \bigcup_{i \in I} Y_i$,
and $f^{-1}(Y_i) = \bigcup X_{ij}$ such that
each morphism $X_{ij} \to Y_i$ is locally of finite type.
\end{enumerate}
\end{lemma}

\begin{proof}
Each of the conditions (2), (3), (4), (5), (6), (7), and (9)
imply condition (8) in a straightforward manner. For example, if
(5) holds, then we can choose a scheme $V$ which is a disjoint
union of affines and a surjective morphism $V \to Y$
(see Properties of Spaces, Lemma
\ref{spaces-properties-lemma-cover-by-union-affines}).
Then $V \times_Y X \to V$ is locally of finite type by (5).
Choose a scheme $U$ and a surjective \'etale morphism
$U \to V \times_Y X$. Then $U \to V$ is locally of finite
type by Lemma \ref{lemma-composition-finite-type}.
Hence (8) is true.

\medskip\noindent
The conditions (1), (7), and (8) are equivalent by definition.

\medskip\noindent
To finish the proof, we show that (1) implies all of the conditions
(2), (3), (4), (5), (6), and (9). For (2) this is immediate.
For (3), (4), (5), and (9) this follows from the fact that being
locally of finite type is preserved under base change, see
Lemma \ref{lemma-base-change-finite-type}.
For (6) we can take $U = X$ and we're done.
\end{proof}

\end{document}
